% Figure 1 caption and surrounding discussion for the quadratic-flow experiment.
% The numerical settings here mirror illustration/hausdorff_experiment.py.

\begin{figure}[H]
  \centering
  \begin{subfigure}[t]{0.335\textwidth}
    \centering
    \includegraphics[width=\linewidth]{fig/figure1_snapshots.png}
    \caption{Reachable-set snapshots.}
    \label{fig:sample_flow_schematic}
  \end{subfigure}
  \hspace{-0.6em}
  \begin{subfigure}[t]{0.335\textwidth}
    \centering
    \includegraphics[width=\linewidth]{fig/figure1_error_vs_time.png}
    \caption{Error versus time.}
    \label{fig:hausdorff_vs_time}
  \end{subfigure}
  \hspace{-0.6em}
  \begin{subfigure}[t]{0.335\textwidth}
    \centering
    \includegraphics[width=\linewidth]{fig/figure1_error_vs_samples.png}
    \caption{Error versus sample size.}
    \label{fig:hausdorff_vs_samples}
  \end{subfigure}
  \caption{Reachable-set approximation for the autonomous non-Lipschitz
  dynamics $\dot y=0$, $\dot x=x^2$.  All initial points are sampled
  uniformly, with the disk shown in red and the star in blue.  Panel (a)
  shows the initial samples and reference boundaries together with three
  propagated snapshots before the finite-time blow-up horizon.  Panel (b)
  shows the mean symmetric Hausdorff distance
  $d_H(S_T,\widehat S_N)$ versus $T$ at $N=1000$.  Panel (c) shows
  the same error versus
  $N\in\{10,30,100,300,1000,3000,10000\}$ at the fixed horizon $T=0.22$ on a
  logarithmic $N$ axis.  The estimator is the convex hull of the propagated
  samples, discretized on a $170\times170$ grid; the reference reachable set is
  represented by $35{,}000$ uniformly sampled points.  Each curve reports the
  mean over 50 trials and the shaded region is the normal-approximation 95\%
  confidence interval for that mean (mean $\pm1.96$ standard errors).  The
  reported symmetric cloud distance is
  $\max\{\max_{s\in S_T}\min_{\hat s\in\widehat S_N}\|s-\hat s\|_2,
  \max_{\hat s\in\widehat S_N}\min_{s\in S_T}\|\hat s-s\|_2\}$.}
  \label{fig:toy_example}
\end{figure}

Figure~\ref{fig:toy_example} isolates the empirical effect of the
non-Lipschitz quadratic flow on sampling-based reachable-set approximation.
As trajectories approach the finite-time blow-up horizon, flow sensitivity and
the resulting Hausdorff approximation error increase sharply.  No theoretical
Lipschitz upper or lower bound is asserted for this example.

At the sample-dependence horizon $T=0.22$, the propagated disk remains convex
and is compatible with the convex-hull estimator, whereas the nonconvex
polygonal star is not.  In particular, the star has zero reach because of its
re-entrant corners, and convexification fills its concavities.  Consequently,
the star curve can retain an irreducible estimator bias as $N$ grows; it should
not be interpreted as evidence of convergence to zero.  Moreover, sufficiently
close to blow-up even the propagated disk can become nonconvex, so late-time
errors can contain convexification bias in addition to sampling error and flow
sensitivity.  A nonconvex support estimator would be required to isolate
zero-error sampling convergence for either nonconvex reachable geometry.
